\documentclass[10pt,a4paper]{amsart}
\usepackage[margin=19mm]{geometry}
\usepackage{amsmath,amssymb,mathtools}
\usepackage[scr=rsfs,cal=euler]{mathalfa}
\usepackage{xfrac}
\usepackage[unicode,hidelinks,bookmarks=false]{hyperref}
\newcommand{\ie}{i.e.\ }
\newcommand{\cf}{cf.\ }
\newcommand{\resp}{respectively\ }
\newcommand{\Subsection}{\subsection}
\providecommand{\nobreakdash}{\nobreak\hbox{-}}
\setlength{\emergencystretch}{2em}
\multlinegap0pt
\usepackage{amssymb}
\usepackage{algorithm}\usepackage{algpseudocode}
\usepackage[shortlabels]{enumitem}
\usepackage[normalem]{ulem}
\usepackage{xcolor}
\newtheorem{lemma}{Lemma}
\title{Perturbation Method in Musielak-Orlicz Sequence Spaces}
\author{Pando Georgiev, Vasil Zhelinski, Boyan Zlatanov}
\date{Source: arXiv:2603.28404v2 (CC BY 4.0)}
\begin{document}
\maketitle

\noindent\emph{Source and licence.} Perturbation Method in Musielak-Orlicz Sequence Spaces. Version arXiv:2603.28404v2 (CC BY 4.0), distributed by arXiv under the Creative Commons Attribution 4.0 International licence, http://creativecommons.org/licenses/by/4.0/, as recorded in the arXiv licence metadata for this exact version and shown on the abstract page https://arxiv.org/abs/2603.28404v2. Full licence notice: \texttt{licenses/arxiv-2603.28404-CC-BY-4.0.txt}.\par\smallskip

\noindent\textit{Editorial scope.} Counted unit: the article's lemma lemma:8A (source0.tex lines 754--777). The article's complete original proof of it is delivered with it (source0.tex lines 780--1015, one \texttt{proof} environment of 236 lines). No mathematics is written, repaired or re-derived: the statement and the proof are the article's own lines, in the article's own notation, and no retained line is substituted or re-flowed.  No further source-local context is needed: every cross-reference of every retained line resolves inside the two delivered spans.  Nothing was omitted from any retained span, so the package carries no omission record.  No retained line contains a citation, so the wrapper carries no bibliography and no bibliography entry.  No retained line contains a figure environment, an image-inclusion command or drawing code, so no image is delivered and the package declares no asset dependency.  Editorial formatting only, in addition: an AMS \texttt{amsart} preamble that restates the article's own declarations and macros used by the retained lines, namely the environments \texttt{lemma} on separate counters, together with the standard packages those lines need.  The retained environments print in this document as Lemma 1 in order of appearance, whereas the article prints the same items with the numbering of its own full document.  The complete native source of the article as distributed in the arXiv e-print for this exact version is bundled unchanged as \texttt{sources/197435/source0.tex}.
\begin{lemma}\label{lemma:8A}
Let $\Phi$ be a non--degenerate Musielak--Orlicz function.  
Suppose there exists a sequence $\{x^{(n)}\}_{n=1}^\infty \subset \ell_\Phi$ such that


\[
I := \inf_{n\in\mathbb N} \|x^{(n)}\|_\Phi > 0,
\]


and


\[
\lim_{n\to\infty} e_\Phi(x^{(n)})
= \lim_{n\to\infty} \sum_{i=1}^\infty \Phi_i\!\left(x^{(n)}_i\right)
= 0.
\tag{9}
\]


Then $h_\Phi \not\cong \ell_\Phi$.  
In particular, there exists $x\in \ell_\Phi$ such that $x\notin h_\Phi$.
\end{lemma}

\begin{proof}
{%\color{ForestGreen}
Since $\|x^{(n)}\|_\Phi \ge I$, by the definition of the Luxemburg norm,


\[
e_\Phi\!\left(\frac{x^{(n)}}{I}\right)
= \sum_{i=1}^\infty \Phi_i\!\left(\frac{x^{(n)}_i}{I}\right)
\ge 1,
\qquad n\in\mathbb N.
\tag{10}
\]



We consider two cases.

\medskip
\noindent
\textbf{Case I.}  
Assume that for every $i\in\mathbb N$,


\[
\lim_{n\to\infty}
\Phi_i\!\left(\frac{x^{(n)}_i}{I}\right) = 0.
\]



\medskip
\textbf{Step 1: Small initial blocks.}
Fix $m\in\mathbb N$.  
Since each term tends to $0$, the finite sum


\[
\sum_{i=1}^m \Phi_i\!\left(\frac{x^{(n)}_i}{I}\right)
\to 0.
\]


Thus there exists $N(m)$ such that for all $n\ge N(m)$,


\[
e_\Phi\!\left(\sum_{i=1}^m \frac{x^{(n)}_i}{I} e_i\right)
\le \frac14.
\tag{11}
\]



\medskip
\textbf{Step 2: Small tails.}
For each $n$, since $e_\Phi(x^{(n)}/I)\ge 1$ and finite,  
there exists $p_n\in\mathbb N$ such that


\[
\sum_{i=p_n+1}^\infty 
\Phi_i\!\left(\frac{x^{(n)}_i}{I}\right)
\le \frac14.
\tag{13}
\]



\medskip
\textbf{Step 3: A large middle block.}
Let $N_1,N_2\in\mathbb N$.  
Set $m=N_2-1$ and choose $n\ge \max\{N_1,\,N(m)\}$.

Then:

- by (11), the block $1,\dots,N_2-1$ contributes at most $1/4$,
- by (13), the tail $p_n+1,\dots,\infty$ contributes at most $1/4$,
- by (10), the whole vector contributes at least $1$.

Hence the middle block satisfies


\[
e_\Phi\!\left(\sum_{i=N_2}^{p_n} \frac{x^{(n)}_i}{I} e_i\right)
\ge \frac12.
\tag{14}
\]



\medskip
\textbf{Step 4: Choosing blocks with small unscaled modular.}
From (9), $e_\Phi(x^{(n)})\to 0$.  
Thus for each $n,N$ there exists $j_{N,n}$ such that


\[
e_\Phi\!\left(\sum_{i=N}^{p_{j_{N,n}}} x^{(j_{N,n})}_i e_i\right)
\le \frac{1}{2n},
\]


while by (14),


\[
e_\Phi\!\left(\sum_{i=N}^{p_{j_{N,n}}}
\frac{x^{(j_{N,n})}_i}{I} e_i\right)
\ge \frac12.
\tag{15}
\]



\medskip
\textbf{Step 5: Diagonal block construction.}
Define $q_1=0$ and recursively


\[
q_{n+1} := p_{\,j_{q_n+1,n}}.
\]


Then $q_1<q_2<q_3<\cdots$.

Define


\[
x^\ast := 
\sum_{n=1}^\infty 
\sum_{i=q_n+1}^{q_{n+1}}
x^{(j_{q_n+1,n})}_i\, e_i.
\]



Since the blocks are disjoint,


\[
e_\Phi(x^\ast)
= \sum_{n=1}^\infty 
e_\Phi\!\left(\sum_{i=q_n+1}^{q_{n+1}}
x^{(j_{q_n+1,n})}_i e_i\right)
\le \sum_{n=1}^\infty \frac{1}{2^n}
< \infty,
\]


so $x^\ast\in \ell_\Phi$.

But also,


\[
e_\Phi\!\left(\frac{x^\ast}{I}\right)
= \sum_{n=1}^\infty 
e_\Phi\!\left(\sum_{i=q_n+1}^{q_{n+1}}
\frac{x^{(j_{q_n+1,n})}_i}{I} e_i\right)
\ge \sum_{n=1}^\infty \frac12
= \infty,
\]


so $x^\ast\notin h_\Phi$.

Thus $\ell_\Phi\neq h_\Phi$ in Case I.

\bigskip
\noindent
\textbf{Case II.}  
Assume Case I fails.  
Then there exists $i_0$ such that


\[
\limsup_{n\to\infty}
\Phi_{i_0}\!\left(\frac{x^{(n)}_{i_0}}{I}\right) > 0.
\]


Choose a subsequence $\{x^{(n_k)}\}$ such that


\[
\Phi_{i_0}\!\left(\frac{x^{(n_k)}_{i_0}}{I}\right)
\to c_0 > 0.
\tag{16}
\]



From (9),


\[
\Phi_{i_0}\!\left(x^{(n_k)}_{i_0}\right)
\le e_\Phi(x^{(n_k)}) \to 0.
\tag{18}
\]



Let $r_k := x^{(n_k)}_{i_0}$.  
Then


\[
\Phi_{i_0}(r_k)\to 0,
\qquad
\Phi_{i_0}\!\left(\frac{r_k}{I}\right)\to c_0>0.
\tag{19}
\]



Since $\Phi_{i_0}$ is non--degenerate,  
$\Phi_{i_0}(t_n)\to 0$ implies $t_n\to 0$.  
Thus $r_k\to 0$, hence $r_k/I\to 0$, and by continuity,


\[
\Phi_{i_0}\!\left(\frac{r_k}{I}\right)\to 0,
\]


contradicting (19).

\bigskip
\noindent
Since Case II is impossible, Case I must occur, and in that case we constructed
$x^\ast\in \ell_\Phi\setminus h_\Phi$.  
Thus $h_\Phi\not\cong \ell_\Phi$.}
\end{proof}

\end{document}
